\documentclass[11pt,a4paper]{article}
\usepackage[T1]{fontenc}
\usepackage[utf8]{inputenc}
\usepackage{amsmath,amssymb,amsthm}
\usepackage[margin=1in]{geometry}
\usepackage[colorlinks=true,linkcolor=blue,urlcolor=blue]{hyperref}
\theoremstyle{plain}
\newtheorem{theorem}{Theorem}
\newtheorem{lemma}[theorem]{Lemma}
\newtheorem{proposition}[theorem]{Proposition}
\newtheorem{corollary}[theorem]{Corollary}
\theoremstyle{definition}
\newtheorem{remark}[theorem]{Remark}
\title{A disproof of the conjectured recurrence for OEIS A076217}
\author{Adrian Perez Fontelles\\ \small Independent researcher}
\date{31 August 2026}
\begin{document}
\maketitle

\begin{abstract}
OEIS A076217 carries a conjectured linear recurrence of order three, contributed by Colin
Barker. We disprove it, and not merely by exhibiting a counterexample: we prove that the
recurrence fails at $n=3^{k}$, $n=3^{k}+1$ and $n=3^{k}+2$ for \emph{every} integer
$k\ge 2$, so the set of indices at which it fails is infinite. The proof is elementary and
self-contained. Everything follows from one fact about the sequence, that $a(n)=1$ exactly
when $n=3^{k}-2$, which is established by an induction driven directly by the entry's own
defining recursion; the values at the six indices surrounding each power of three are then
computed in closed form, and the failure is read off. The entry records the observation
that the recurrence ``seems to fail at $n=$ powers of $3$''; what is proved here is that it
does, always.
\end{abstract}

\noindent\small 2020 Mathematics Subject Classification. 11B37, 11B83.\normalsize

\section{The sequence and the conjecture}

OEIS A076217 is defined by its name, which is the recursion itself:
\begin{equation}\label{eq:def}
a(1)=1,\qquad a(n)=a(n-1)+n\cdot\operatorname{sign}\bigl(n-a(n-1)\bigr)\quad(n\ge2),
\end{equation}
where $\operatorname{sign}(x)$ is $1$, $0$ or $-1$ according as $x$ is positive, zero or
negative. It has offset $1$ and begins
\[
a(1),\dots,a(12)\;=\;1,\ 3,\ 3,\ 7,\ 2,\ 8,\ 1,\ 9,\ 9,\ 19,\ 8,\ 20,\ \dots
\]
Nothing else is taken from the entry: \eqref{eq:def} is the definition, so every value used
below is computed from it.

The entry carries the following comment:
\begin{quote}\small
Conjecture: a(n) = -a(n-1)+a(n-2)+a(n-3) for n>5. G.f.: $-x(27x^{28}+54x^{27}+27x^{26}
+9x^{10}+18x^{9}+9x^{8}+3x^{4}+6x^{3}+5x^{2}+4x+1)/((x-1)(x+1)^{2})$. --- \emph{Colin
Barker}, Feb 25 2013
\end{quote}
As of the ``Last modified'' line on the live entry (2025-03-21, revision 33) the statement
is still recorded as a conjecture. The entry also carries a later remark that the
recurrence ``seems to fail at $n=$ powers of 3, and the 2 successive terms''. That is an
observation about the terms in range; the theorem below proves it for all $k$, and the
conjectured generating function falls with the recurrence, since a rational function with
denominator $(x-1)(x+1)^{2}$ forces
\begin{equation}\label{eq:conj}
a(n)+a(n-1)-a(n-2)-a(n-3)=0
\end{equation}
for all $n$ past the degree of its numerator.

\section{Where the sequence takes the value one}

Write $p=3^{k}$ throughout. The whole argument rests on one proposition.

\begin{proposition}\label{prop:ones}
For every $k\ge1$, $a(3^{k}-2)=1$.
\end{proposition}

The proof is an induction, and the induction step also produces the closed forms we need.

\begin{lemma}\label{lem:block}
Let $p\ge3$ and suppose $a(p-2)=1$. Then
\[
a(p-1)=p,\qquad a(p)=p,
\]
and for every $i$ with $0\le i\le p-2$,
\[
a(p+1+2i)=2p+1+i,\qquad a(p+2+2i)=p-1-i .
\]
\end{lemma}

\begin{proof}
All four statements are applications of \eqref{eq:def}, and in each one the only thing to
check is the sign of $n-a(n-1)$.

At $n=p-1$ the previous value is $a(p-2)=1$ and $n-a(n-1)=p-2>0$, so the sign is $+1$ and
$a(p-1)=1+(p-1)=p$.

At $n=p$ the previous value is $a(p-1)=p$, so $n-a(n-1)=0$, the sign is $0$, and
$a(p)=a(p-1)=p$.

Now argue by induction on $i$. For the odd-offset index $n=p+1+2i$ the previous value is
$a(p+2i)$, which is $a(p)=p$ when $i=0$ and $p-1-(i-1)=p-i$ when $i\ge1$; in both cases it
equals $p-i$ for $i\ge 1$ and $p$ for $i=0$, and in both cases
\[
n-a(n-1)\;=\;(p+1+2i)-(p-i)\;=\;1+3i\;>\;0 \qquad(i\ge1),
\]
while for $i=0$ it is $(p+1)-p=1>0$. So the sign is $+1$ and
\[
a(p+1+2i)=(p-i)+(p+1+2i)=2p+1+i
\]
for $i\ge1$, and $a(p+1)=p+(p+1)=2p+1$ for $i=0$, which is the same formula.

For the even-offset index $n=p+2+2i$ the previous value is $a(p+1+2i)=2p+1+i$, and
\[
a(n-1)-n\;=\;(2p+1+i)-(p+2+2i)\;=\;p-1-i,
\]
which is positive exactly when $i\le p-2$. So for those $i$ the sign is $-1$ and
\[
a(p+2+2i)=(2p+1+i)-(p+2+2i)=p-1-i,
\]
as claimed. This completes the induction on $i$ and the lemma.
\end{proof}

\begin{proof}[Proof of Proposition~\ref{prop:ones}]
For $k=1$ we have $3^{1}-2=1$ and $a(1)=1$ by definition. Suppose $a(3^{k}-2)=1$ for some
$k\ge1$ and put $p=3^{k}\ge3$. Lemma~\ref{lem:block} applies. Take $i=p-2$, which is within
the permitted range $0\le i\le p-2$; then
\[
p+2+2i\;=\;p+2+2(p-2)\;=\;3p-2\;=\;3^{k+1}-2,
\]
and the lemma gives
\[
a\bigl(3^{k+1}-2\bigr)\;=\;p-1-i\;=\;p-1-(p-2)\;=\;1 .
\]
That is the statement for $k+1$.
\end{proof}

\begin{corollary}\label{cor:vals}
For every $k\ge1$, writing $p=3^{k}$,
\[
a(p-3)=p-1,\quad a(p-2)=1,\quad a(p-1)=p,\quad a(p)=p,\quad a(p+1)=2p+1,\quad a(p+2)=p-1 .
\]
\end{corollary}

\begin{proof}
The values at $p-2$, $p-1$, $p$, $p+1$ and $p+2$ are Proposition~\ref{prop:ones} together
with Lemma~\ref{lem:block} at $i=0$.

For $a(p-3)$, apply Lemma~\ref{lem:block} at the previous power. Put $q=3^{k-1}$, so
$p=3q$ and $p-3=3q-3$. The lemma describes the indices after $q$ in two families,
$q+1+2i$ and $q+2+2i$. Solving $q+2+2i=3q-3$ gives $2i=2q-5$, which has no integer
solution, so $p-3$ is not of the second kind. Solving $q+1+2i=3q-3$ gives
\[
i\;=\;q-2,
\]
which lies in the permitted range $0\le i\le q-2$, and the first family gives
\[
a(p-3)\;=\;2q+1+i\;=\;2q+1+(q-2)\;=\;3q-1\;=\;p-1 .
\]
This is legitimate for $k\ge2$, since then $q\ge3$ and Lemma~\ref{lem:block} applies at
$q$; and $i=q-2$ is the same index at which the lemma yields $a(3q-2)=1$, so both values
come from a single application.
\end{proof}

\section{The recurrence fails at every power of three}

\begin{theorem}
For every integer $k\ge2$ the conjectured recurrence \eqref{eq:conj} fails at $n=3^{k}$, at
$n=3^{k}+1$ and at $n=3^{k}+2$. In particular it fails at infinitely many indices, so it
does not hold for all $n>5$, and the conjectured generating function is likewise incorrect.
\end{theorem}

\begin{proof}
Put $p=3^{k}$ and use Corollary~\ref{cor:vals}.

At $n=p$ the right-hand side of the conjecture is
\[
-a(p-1)+a(p-2)+a(p-3)\;=\;-p+1+(p-1)\;=\;0,
\]
whereas $a(p)=p$. Since $p=3^{k}\ge9>0$, the two differ.

At $n=p+1$ it is
\[
-a(p)+a(p-1)+a(p-2)\;=\;-p+p+1\;=\;1,
\]
whereas $a(p+1)=2p+1\ge19$.

At $n=p+2$ it is
\[
-a(p+1)+a(p)+a(p-1)\;=\;-(2p+1)+p+p\;=\;-1,
\]
whereas $a(p+2)=p-1\ge8$.

All three indices exceed $5$, so each is inside the range the conjecture claims. As $k$
ranges over the integers $\ge2$ these give infinitely many failures.
\end{proof}

\section{Verification}

Two checks were made, independent of the algebra above.

First, the recursion \eqref{eq:def} was iterated in exact integer arithmetic to $n=60000$
and the resulting values compared with the entry's DATA field and with its b-file, which
carries $10000$ terms; they agree at every index.

Second, the conjectured recurrence was evaluated on those values at every index from $6$ to
$4000$. It fails at exactly the indices
\[
9,10,11,\quad 27,28,29,\quad 81,82,83,\quad 243,244,245,\quad 729,730,731,\quad
2187,2188,2189,
\]
that is, at $3^{k},3^{k}+1,3^{k}+2$ for $k=2,\dots,7$ and nowhere else, matching the
theorem exactly. The closed forms of Corollary~\ref{cor:vals} were checked at every
$k\le 9$ and the block formulas of Lemma~\ref{lem:block} at every admissible $i$ in that
range.

\section{Concluding remarks}

The conjecture is not close to true: it fails on a set of indices of the shape
$3^{k}+\{0,1,2\}$, and between consecutive powers of three it holds only because the
sequence is locked into the simple alternating pattern of Lemma~\ref{lem:block}, which the
recurrence happens to reproduce. The pattern is broken precisely at the moment the downward
branch reaches the value $1$, which by Proposition~\ref{prop:ones} happens exactly at
$n=3^{k}-2$.

\begin{thebibliography}{9}
\bibitem{oeis} OEIS Foundation Inc., \emph{The On-Line Encyclopedia of Integer Sequences},
entry \href{https://oeis.org/A076217}{A076217}, 2025.
\end{thebibliography}

\end{document}
